% GENERATED FILE. Edit canonical lessons, then run npm run generate:books.
% canonical-source-hash: fnv1a64:36c2e7c066683d0a
% canonical-lessons: AR-C423-bawwaba, AR-C423-maqad, AR-C423-hajaza, AR-C423-suud, AR-C423-madraj

\chapter{At the gate}
\label{ch:ar-at-the-gate}
\hlchaptermodality{423}

\begin{chapteropening}
\textbf{By the end of this chapter:} \emph{I can say a gate, a seat, to reserve, boarding and a runway.}

Complete the last lesson of chapter 423: I can say a gate, a seat, to reserve, boarding and a runway.
\end{chapteropening}

\section[\texorpdfstring{\ar{بوابة} (bawwāba)}{بوابة (bawwāba)}]{\ar{بوابة} (bawwāba) — a gate; a gateway}

\label{lesson:AR-C423-bawwaba}



\begin{quote}
Before the new one: say the Arabic for a sentence, then the Arabic for a barber.
\end{quote}

\subsection*{You'll want to know: \ar{بوابة}}
\textbf{\ar{بوابة}} — \emph{bawwāba} — \textquotedblleft{}a gate; a gateway\textquotedblright{}.

\begin{quote}
\ar{انتظر} \ar{المسافرون} \ar{عند} \ar{البوابة}.
\end{quote}

\begin{quote}
\emph{intaẓara al-musāfirūn ʿinda al-bawwāba.}
\end{quote}

\begin{quote}
The travellers waited at the gate.
\end{quote}

\begin{cousinweb}
\textbf{\ar{بوابة}} (\emph{bawwāba}) is built on the native Arabic root \emph{b-w-b}, \textquotedblleft{}a door\textquotedblright{} (Lane). It shares its root with \emph{bāb}, \textquotedblleft{}a door\textquotedblright{}; Lane gives \emph{bawwāb}, \textquotedblleft{}a door-keeper\textquotedblright{}, on the same root, and modern Arabic uses \emph{bawwāba} for a gate or a gateway.
\end{cousinweb}

\begin{grammarlens}[title={What you've built}]
One more word for getting around.
\end{grammarlens}

\subsection*{Your turn}
\begin{itemize}
\raggedright
  \item \emph{Say it:} \emph{bawwāba}
  \item \emph{Say it:} \emph{bawwāba}, once more
  \item \emph{Say it:} \emph{intaẓara al-musāfirūn ʿinda al-bawwāba.} — The travellers waited at the gate.
  \item \emph{Recall:} say the Arabic for a sentence, then the Arabic for a barber, then say \emph{bawwāba} again
\end{itemize}

\begin{tcolorbox}[breakable,colback=teal!4,colframe=teal!35!black,title={Before you move on}]
What does \ar{بوابة} mean? (\textquotedblleft{}A gate; a gateway\textquotedblright{}.) Say it once more.
\end{tcolorbox}
\section[\texorpdfstring{\ar{مقعد} (maqʿad)}{مقعد (maqʿad)}]{\ar{مقعد} (maqʿad) — a seat}

\label{lesson:AR-C423-maqad}



\begin{quote}
Before the new one: say the Arabic for a barber, then the Arabic for a gate.
\end{quote}

\subsection*{You'll want to know: \ar{مقعد}}
\textbf{\ar{مقعد}} — \emph{maqʿad} — \textquotedblleft{}a seat\textquotedblright{}.

\begin{quote}
\ar{جلست} \ar{في} \ar{مقعد} \ar{بجانب} \ar{النافذة}.
\end{quote}

\begin{quote}
\emph{jalastu fī maqʿad bi-jānib an-nāfidha.}
\end{quote}

\begin{quote}
I sat in a seat beside the window.
\end{quote}

Plural: \ar{مقاعد}, \emph{maqāʿid}.

The pattern \emph{mafʿal} names a place: \emph{maqʿad} is the place where one sits.

\begin{cousinweb}
\textbf{\ar{مقعد}} (\emph{maqʿad}) is built on the native Arabic root \emph{q-ʿ-d}, \textquotedblleft{}to sit down\textquotedblright{} (Lane).
\end{cousinweb}

\begin{grammarlens}[title={What you've built}]
One more word for getting around.
\end{grammarlens}

\subsection*{Your turn}
\begin{itemize}
\raggedright
  \item \emph{Say it:} \emph{maqʿad}
  \item \emph{Say it:} \emph{maqʿad}, once more
  \item \emph{Say it:} \emph{jalastu fī maqʿad bi-jānib an-nāfidha.} — I sat in a seat beside the window.
  \item \emph{Recall:} say the Arabic for a barber, then the Arabic for a gate, then say \emph{maqʿad} again
\end{itemize}

\begin{tcolorbox}[breakable,colback=teal!4,colframe=teal!35!black,title={Before you move on}]
What does \ar{مقعد} mean? (\textquotedblleft{}A seat\textquotedblright{}.) Say it once more.
\end{tcolorbox}
\section[\texorpdfstring{\ar{حجز} (ḥajaza)}{حجز (ḥajaza)}]{\ar{حجز} (ḥajaza) — to reserve; to book}

\label{lesson:AR-C423-hajaza}



\begin{quote}
Before the new one: say the Arabic for a gate, then the Arabic for a seat.
\end{quote}

\subsection*{You'll want to know: \ar{حجز}}
\textbf{\ar{حجز}} — \emph{ḥajaza} — \textquotedblleft{}to reserve; to book\textquotedblright{}.

\begin{quote}
\ar{حجزت} \ar{مقعدا} \ar{في} \ar{القطار}.
\end{quote}

\begin{quote}
\emph{ḥajaztu maqʿadan fī al-qiṭār.}
\end{quote}

\begin{quote}
I booked a seat on the train.
\end{quote}

Present: \ar{يحجز}, \emph{yaḥjiz}.

Lane gives the present both as \emph{yaḥjizu} and \emph{yaḥjuzu}. The noun \emph{ḥajz} is a booking.

\begin{cousinweb}
\textbf{\ar{حجز}} (\emph{ḥajaza}) is built on the native Arabic root \emph{ḥ-j-z}, \textquotedblleft{}to hold back, to withhold\textquotedblright{} (Lane). A booking holds a place back for someone. It is the plain verb.
\end{cousinweb}

\begin{grammarlens}[title={What you've built}]
One more word for getting around.
\end{grammarlens}

\subsection*{Your turn}
\begin{itemize}
\raggedright
  \item \emph{Say it:} \emph{ḥajaza}
  \item \emph{Say it:} \emph{ḥajaza}, once more
  \item \emph{Say it:} \emph{ḥajaztu maqʿadan fī al-qiṭār.} — I booked a seat on the train.
  \item \emph{Recall:} say the Arabic for a gate, then the Arabic for a seat, then say \emph{ḥajaza} again
\end{itemize}

\begin{tcolorbox}[breakable,colback=teal!4,colframe=teal!35!black,title={Before you move on}]
What does \ar{حجز} mean? (\textquotedblleft{}To reserve; to book\textquotedblright{}.) Say it once more.
\end{tcolorbox}
\section[\texorpdfstring{\ar{صعود} (ṣuʿūd)}{صعود (ṣuʿūd)}]{\ar{صعود} (ṣuʿūd) — boarding; going up}

\label{lesson:AR-C423-suud}



\begin{quote}
Before the new one: say the Arabic for a seat, then the Arabic for to reserve.
\end{quote}

\subsection*{You'll want to know: \ar{صعود}}
\textbf{\ar{صعود}} — \emph{ṣuʿūd} — \textquotedblleft{}boarding; going up\textquotedblright{}.

\begin{quote}
\ar{بدأ} \ar{الصعود} \ar{بعد} \ar{التأخير}.
\end{quote}

\begin{quote}
\emph{badaʾa aṣ-ṣuʿūd baʿda at-taʾkhīr.}
\end{quote}

\begin{quote}
Boarding began after the delay.
\end{quote}

\begin{cousinweb}
\textbf{\ar{صعود}} (\emph{ṣuʿūd}) is built on the native Arabic root \emph{ṣ-ʿ-d}, \textquotedblleft{}to go up, to ascend\textquotedblright{} (Lane). It is the verbal noun of \emph{ṣaʿida}, \textquotedblleft{}to go up\textquotedblright{}.
\end{cousinweb}

\begin{grammarlens}[title={What you've built}]
One more word for getting around.
\end{grammarlens}

\subsection*{Your turn}
\begin{itemize}
\raggedright
  \item \emph{Say it:} \emph{ṣuʿūd}
  \item \emph{Say it:} \emph{ṣuʿūd}, once more
  \item \emph{Say it:} \emph{badaʾa aṣ-ṣuʿūd baʿda at-taʾkhīr.} — Boarding began after the delay.
  \item \emph{Recall:} say the Arabic for a seat, then the Arabic for to reserve, then say \emph{ṣuʿūd} again
\end{itemize}

\begin{tcolorbox}[breakable,colback=teal!4,colframe=teal!35!black,title={Before you move on}]
What does \ar{صعود} mean? (\textquotedblleft{}Boarding; going up\textquotedblright{}.) Say it once more.
\end{tcolorbox}
\section[\texorpdfstring{\ar{مدرج} (madraj)}{مدرج (madraj)}]{\ar{مدرج} (madraj) — a runway}

\label{lesson:AR-C423-madraj}



\begin{quote}
Before the new one: say the Arabic for to reserve, then the Arabic for boarding.
\end{quote}

\subsection*{You'll want to know: \ar{مدرج}}
\textbf{\ar{مدرج}} — \emph{madraj} — \textquotedblleft{}a runway\textquotedblright{}.

\begin{quote}
\ar{كانت} \ar{الرحلة} \ar{تنتظر} \ar{على} \ar{المدرج}.
\end{quote}

\begin{quote}
\emph{kānat ar-riḥla tantaẓiru ʿalā al-madraj.}
\end{quote}

\begin{quote}
The flight was waiting on the runway.
\end{quote}

Plural: \ar{مدارج}, \emph{madārij}.

\begin{cousinweb}
\textbf{\ar{مدرج}} (\emph{madraj}) is built on the native Arabic root \emph{d-r-j}, \textquotedblleft{}to go step by step\textquotedblright{} (Lane). Lane gives \emph{madraj} as a way or route one goes along; modern Arabic gave the word the runway.
\end{cousinweb}

\begin{grammarlens}[title={What you've built}]
One more word for getting around.
\end{grammarlens}

\subsection*{Your turn}
\begin{itemize}
\raggedright
  \item \emph{Say it:} \emph{madraj}
  \item \emph{Say it:} \emph{madraj}, once more
  \item \emph{Say it:} \emph{kānat ar-riḥla tantaẓiru ʿalā al-madraj.} — The flight was waiting on the runway.
  \item \emph{Recall:} say the Arabic for to reserve, then the Arabic for boarding, then say \emph{madraj} again
\end{itemize}

\begin{tcolorbox}[breakable,colback=teal!4,colframe=teal!35!black,title={Before you move on}]
What does \ar{مدرج} mean? (\textquotedblleft{}A runway\textquotedblright{}.) Say it once more.
\end{tcolorbox}
